% 本文件由 LaTeX 排版 Agent 根据有效 translation.json 自动生成。
% author=Dionysius the Great; work_id=0612; title=杂著

\chapter{杂著}

% structure: p0001 source=h1 target=omitted reason=page-title-already-rendered-by-parent

% structure: p0002 source=h2 target=section reason=relative-heading-rank; base=h2

\section{I. 取自《论预许》两卷书}

1. 但他们拿出奈波斯的一篇作品，极力坚持，似乎它不可辩驳地证明了基督将在地上有一个（暂时的）统治。我必须说，在其它许多方面，我赞同奈波斯的意见，并因他的信德、勤勉以及在圣经上的潜心研读而爱他，也因他在圣咏方面的巨大努力（甚至现在许多弟兄仍因此欣喜）而爱他。我更深深敬重此人，因为他已先我们安息。然而，真理应被珍视并敬拜于一切之上。的确，凡正确言论都应当不吝赞许和认可，但同样也应当审查和纠正任何可能被视为不健全的写作。如果他当时亲自在场，并以口头陈述他的见解，那么在没有写作的情况下一起讨论问题，并试图说服反对者、通过问答引导他们，就已足够。但作品已经发表，而且对某些人似乎极具说服力；并且毫无疑问，有些教师认为律法和先知无关紧要，拒绝跟随福音，贬低宗徒的书信，并对这篇作品所宣扬的教理做出重大许诺，好像它是某种伟大而隐藏的奥秘；与此同时，他们又不允许我们较为单纯的弟兄对主的荣耀显现与祂的真实神性，以及我们自己的死而复活、被聚集到祂面前、与祂相似有任何崇高和卓越的观念，反而竭力引导他们盼望琐碎可朽坏的事物，只如我们如今在天主国内所见的那样。既然如此，我们就有必要如同奈波斯弟兄在场一样，与他讨论这个问题。\par

2. \Emph{在讨论了一些其他问题之后，他补充了以下声明：}——当时我身在阿尔西诺伊特州——正如你们所知，这一教理在那里早已流传，并引起如此分裂，以致整个教会发生了分裂和背弃——我召集了村庄里的长老、教师中弟兄们，以及那些愿意出席的弟兄们也都在场。我劝勉他们公开调查那项教义。因此，当他们把这本书带到我们面前，好像它是一种武器或不可攻破的堡垒时，我和他们从早到晚连续坐了三天，试图纠正他们对作品中所提主题的认识。同时，我非常欣慰地观察到弟兄们的坚定、对真理的热爱、他们的温顺和智慧，因为我们以有序的方法和温和的精神处理问题、疑难和让步。因为我们注意不固执地、嫉妒地坚持已经接受的意见，即使它们看起来可能是正确的。我们也没有回避别人提出的反对意见；而是尽可能努力紧扣主题，确立相关的论点。同样，如果理智说服了我们，我们也不耻于改变意见并承认事实；反而凭着良心，以完全的真诚、向天主敞开的心怀，接受了一切能够通过圣经的证明和教导而确立的内容。最后，这个教义的作者和倡导者，名叫科拉奇翁，在所有在场的弟兄面前承认了他的立场，并向我们保证他不再坚持自己的意见，不再讨论、提及或教导它，因为他已完全被反对者的论证所说服。其余在场的弟兄们也因这次会议以及所有人表现出的和解精神与和谐而感到欣慰。\par

3. \Emph{然后，稍后，他论及\WorkTitle{若望默示录}如下}：——如今有些在我们之前的人弃绝了这卷书，全然否认它，逐章批评，力图证明它毫无意义和理智。他们还声称它的标题是虚假的，因为他们否认若望是作者。不仅如此，他们还认为它根本不是什么启示，因为它是被如此粗俗而浓密的愚昧之幔所遮盖。因此他们断言，没有一个宗徒，甚至没有一个圣人，也没有任何属于教会的人可能是其作者；而是克陵托以及由他所创立并以他的名字命名的克陵托异端派，渴望为他自己所提出的虚构作品附上一个伟大名字的权威，便将该标题冠在书前。因为克陵托所灌输的教义是：基督将有一个地上的王国；而他本人是一个沉溺于肉体享乐、全然属肉体性情的人，就幻想那个王国将由他心中所热衷的那类满足构成——即，在肚腹的享乐和肚腹以下的事情，也就是说，在饮食、嫁娶，以及他以为能以更优雅的方式放纵私欲的其他事物，如节庆、祭献和宰杀祭品。至于我，却不敢弃绝这卷书，因为有许多弟兄珍视它。然而，既然形成了它是一部超出我理解能力的作品的概念，我视它为其所涉各个主题中蕴含着一种隐秘而奇妙的智慧。因为我虽不能领会，仍怀疑其字句之下有更深的意义。我不以自己的理性为标准来衡量和判断它的表达，而是更多迁就于信德，只把它们视为过于高深而无法理解；我并不立刻拒绝我所不明白的，反而更加惊叹，因为我未能洞悉其意。\par

4. \Emph{此后，他查考整部默示录；并证实它绝不可能按照粗浅的字面意义来理解，于是继续写道}：——当先知如今已经完成，可以说，全部预言时，他宣告那些遵守它的人是有福的，并将自己也列入其中：\Quote{有福的，}他说，\BibleQuote{是遵守这书上预言的话的人；我\Emph{若望}就是看见并听见这些事的。}\BibleRef{默 22:7} 因此，此人被称为若望，且这是某位若望的著作，我并不否认。我还进一步承认，它也是某位圣洁、受默感的人的著作。但我不能轻易承认这位就是宗徒，载伯德的儿子，雅各伯的兄弟，与写了\BibleBook{若望}{福音}以及\BibleBook{若望}{一书}的是同一个人。但从两者的风格、表达形式以及全书的布局和措辞，我得出结论：作者并非他。因为这位福音作者在其他地方从未附上自己的名字，不论在福音书中还是在书信中，他从未表明自己。\par

\Emph{不久之后，他又补充道}：——而且若望从未给我们名字，无论是作为他自己的直接称呼（第一人称），还是作为另一人（第三人称）。但默示录的作者在开头就表明自己，因为他说：\BibleQuote{耶稣基督的启示，就是天主赐给他，叫他将必要快成的事指示他的众仆人；他就差遣使者晓谕他的仆人若望。若望便将天主的道和耶稣基督的见证，凡自己所看见的都证明出来。}\BibleRef{默 1:1} 然后他又写了一封信，其中说：\Quote{若望写信给亚细亚的七个教会：愿恩宠与平安赐与你们。} 相反，那位福音作者甚至没有在\BibleBook{若望}{一书}前署名；而是毫不绕弯子，直接以神圣启示本身的奥妙开始，这样说：\BibleQuote{论到那从起初就有的，我们所听见，所看见的。}\BibleRef{若一 1:1} 并且基于这样的启示，主曾称伯多禄为有福的，他说：\BibleQuote{西满巴约纳，你是有福的！因为不是肉和血启示了你，而是我在天之父。}\BibleRef{玛 16:17} 此外，在归属为宗徒若望的\BibleBook{若望}{二书}和\BibleBook{若望}{三书}中，虽然它们确实简短，但若望并没有以名字出现在我们面前；我们只发现匿名的写作：\Quote{长老。} 与此相反，这另一位作者甚至认为只称自己一次然后继续叙述并不足够；他再次提到他的名字，说：\BibleQuote{我若望，你们的兄弟，在耶稣基督的患难、王国和忍耐里有分，为天主的道并为耶稣基督的见证，曾在那名叫帕特摩的海岛上。}\BibleRef{默 1:9} 同样，在结尾处，他这样说：\BibleQuote{那些遵守这书上预言的话的人是有福的；我\Emph{若望}就是看见并听见这些事的。}\BibleRef{默 22:7} 那么，我们应当相信，写这些事的是一个若望，因为他自己告诉我们。\par

5. 然而，这位若望究竟是谁，并不确定。因为他没有像在福音中常做的那样，说自己是主所爱的门徒，或是靠在主怀中的那一位，或是雅各伯的兄弟，或是有幸看见并听到主的那一位。而且，如果他有意明确表明自己，他一定会给我们一些这样的提示。但关于这一切，他什么也没有提供；他只称自己为我们的兄弟和同伴，是耶稣的见证人，并有幸看见并听到这些启示。我也认为，与宗徒若望同名的人有很多，他们出于对他的爱、对他的钦佩和效仿，以及渴望像他那样被主所爱，也采用了同样的称呼，正如我们看到许多信友的孩子被称为保禄和伯多禄一样。此外，在《宗徒大事录》中还提到另一位若望，姓氏为马尔谷，巴尔纳伯和保禄收他为同伴，关于他又说：\Quote{他们也有若望作助手。}\BibleRef{宗 13:5} 但这人是否就是写《默示录》的那一位，我不能说。因为经上并没有记载他和他们一同进入亚细亚。但作者说：\Quote{保禄和他的同伴从帕佛起航，来到旁非里雅的培尔革；若望就离开他们，回了耶路撒冷。}\BibleRef{宗 13:13} 因此，我认为那是亚细亚的另外一人。因为据说在厄弗所有两座坟墓，每座都刻着若望的名字。\par

6. 从思想、表达及其排列来看，可以相当合理地推测这位若望与那一位不同。因为福音与书信彼此一致，并且都以同样的方式开头。前者以\Quote{在起初已有圣言}开始，而后者以\Quote{论到那从起初就有的}开始。前者说：\Quote{圣言成了血肉，居住在我们中间；我们看见了祂的光荣，正如圣父独生者的光荣。}\BibleRef{若 1:14} 后者则稍加变动说了同样的话：\Quote{论到那从起初就有的、我们所听见、所看见、亲眼看过、亲手摸过的生命的圣言——这生命已显现出来。}这些是作为序言引入的，并且正如他在后面所表明的，是为了反对那些否认主曾降来肉身的人。为此，他也谨慎地加上了这些话：\Quote{我们将所见所闻的传报给你们，为使你们也与我们相通；而我们的相通是与圣父和祂的子耶稣基督相通的。}这样，他坚持自己的主题，没有前后矛盾地偏离，而是以相同的标题和相同的措辞贯穿所有内容，我们将略举其中一些。细心的读者会发现，两者中频繁出现\Quote{生命}、\Quote{光明}等词句；还有诸如\Emph{远离黑暗、持守真理、恩宠、喜乐、主的血肉、审判、罪的赦免、天主对我们的爱、我们彼此相爱的命令}；以及我们应遵守所有诫命、对世界、魔鬼、假基督的定罪、圣神的恩许、天主的义子、我们所需的一切信德、\Emph{圣父与圣子}，在各处如此称呼。总之，从整体来看，福音和书信显然具有相同写作风格的特征。然而，《默示录》与此完全不同，并且与它截然有别；我几乎可以说它甚至不相近，也不接界。它与其他这些书卷没有一个共同音节。不仅如此，书信——因为我不提福音——丝毫没有提及或表现出对《默示录》的认识，而《默示录》也同样没有提到书信。然而，保禄在他的书信中暗示了他的启示；但这些启示他并没有单独写下来。\par

7. 此外，根据文辞的差异，有可能证明福音和书信与《默示录》之间的区别。因为前者不仅在希腊语方面没有实际错误，而且在表达、推理以及整个风格结构上都非常优雅。它们远非在措辞上带有任何粗野、语法错误或任何粗俗用语。因为可以推测，作者拥有两种话语的恩赐，主赋予了他这两种能力，即知识和表达的能力。至于后者的作者看见了启示，并领受了知识和预言，我不否认。只是我察觉到他的方言和语言并非精确的希腊文类型，他使用了粗野的习语，有些地方还有语法错误。然而，我们目前无需追究这些。我不希望任何人以为我说这些话是出于嘲笑；我这样做只是为了纠正这些著作之间存在差异的问题。\par

% structure: p0011 source=h2 target=section reason=relative-heading-rank; base=h2

\section{II. 摘自《论自然》}

% structure: p0012 source=h3 target=subsection reason=relative-heading-rank; base=h2

\subsection{一、反对伊壁鸠鲁学派否认天主眷顾并以原子构成宇宙者}

宇宙是一个融贯的整体吗？——按我们自己的判断，也按最睿智的希腊哲学家如柏拉图、毕达哥拉斯、斯多葛派和赫拉克利特的判断，它似乎是如此。或者它是一个二元体，像某些人可能猜测的那样？或者它实际上是多元且无限的，像某些其他人所认为的那样？他们以各种狂想的猜测和幻想的术语用法，试图分割并分解宇宙的基本质料，并声称它是无限、无始、且不受天主眷顾支配的。因为有些人将某些不可毁灭的、最微小的、被认为数量无限的物体命名为\OriginalTerm{原子}{atoms}，并假定存在某个无限广袤的虚空空间，声称这些原子在虚空中偶然地运动，完全随机地彼此碰撞，无规则地旋转，并以多种形式相互混合，彼此结合，从而逐渐形成这个世界及其中的一切物体；甚至，它们构造了无限多的世界。这是伊壁鸠鲁和德谟克利特的观点；只是他们在一点上有分歧：前者假定这些原子都是最微小的，因而不可察觉，而德谟克利特则认为其中也有一些非常大的原子。但两人都认为这些原子确实存在，并因其不可分解的坚固性而得名。还有另外一些人将某些不可分割成部分的物体命名为原子，而这些物体本身就是宇宙的组成部分，万物由它们在未分割状态下构成，又分解回归于它们。据说狄奥多洛斯是给它们命名为不可分割物体的那个人。但据说赫拉克利德给了它们另一个名称，称为\Quote{重量；}；而阿斯克勒庇阿得斯医生也由此获得这一名称。\newline{}\par

% structure: p0014 source=h3 target=subsection reason=relative-heading-rank; base=h2

\subsection{二、基于常见人类类比对这一学说的驳斥。\newline{}}

我们怎能容忍这些人，他们断言所有那些智慧的、因而也是高贵的宇宙构造，都只是普遍偶然的作品？我说的是那些物体，其中每一个单独被造时，以及整个系统，都被祂的命令使之存在的那位视为美好的。因为经上记载：\BibleQuote{天主看了祂所造的一切，认为样样都很好。}\BibleRef{创 1:31}但这些人确实没有反思甚至小型常见事物的类比，这些事物随时可被他们观察到，从中他们可以学到：任何有用且适于服务的物体，都不是无设计或偶然产生的，而是由手工技艺制作，并被设计成满足其适当用途。而当物体失去用途、变得无用时，它也开始无定地破裂，并以各种随意无规则的方式分解、消散其材料，就像最初巧妙构造它的智慧不再控制和支持它。例如，一件外衣不能没有织工而制成，仿佛经线能自行对齐、纬线能自行缠绕；而一旦穿旧，破布就被丢弃。同样，当建造房屋或城市时，石头不会自行放置，仿佛有些自动落在基础上，其他的自己升到各层，而是建造者精心将加工好的石头摆放在适当位置；而如果结构一旦坍塌，石头就分离、坠落、散落各处。同样，造船时，龙骨不会自行铺设，桅杆不会自行竖立在中心，其他所有木料也不会偶然地、自行地各就各位。所谓车上的百根横木也不会自行落到它们各自碰上的空位。相反，木匠在合适的时间以正确的方式组装材料。如果船出海遇难，或车在陆上行驶时撞碎，它们的木料就会散落各处——前者被波浪，后者被撞击的暴力。同样，我们完全可以对这些人说：他们那些闲置、未经加工、无用的原子，是徒然引入的。因此，让他们自己去寻求看见不可见之物，构想不可思议之物吧；不像那一位，他这样向天主承认，只有天主亲自向他显示了这类事物：\Quote{我的眼睛看见了您的作为，那时它们尚未完成。}然而，当他们现在断言所有那些优美高雅的东西（他们宣称这些东西是由原子精细织成的）是被这些没有智慧或知觉的物体自发制造出来时，谁能忍受听他们用这样的词语谈论那些无规则的原子呢？就连自行编织其特有技艺的蜘蛛，也比这些原子更有智慧。\newline{}\par

% structure: p0016 source=h3 target=subsection reason=relative-heading-rank; base=h2

\subsection{三、基于宇宙构造的驳斥。\newline{}}

或者，谁能忍受听到这样的主张：这个由天地构成的伟大居所，因其规模与蕴含的智慧之丰富而被称为\Quote{\OriginalTerm{宇宙}{Cosmos}}，竟是由那些无序而丑陋地运行的原子建立起来的，并且那种无序状态竟演变成了这个真正的宇宙、秩序？或者，谁能相信那些规则的运动和轨迹是某种无规则的冲动的产物？或者，谁能允许天体之间存在的完美和谐，是从缺乏和谐与协调的工具中产生出来的？再者，如果万物只有同一实体，所有事物都具有同一不朽的本性——据他们所言，唯一的差异在于大小和形状——那么，怎么会有一些神圣而完美的、永恒的（如他们所说）或持久的（如有人喜欢表达的）物体呢？其中有些是可见的，有些是不可见的——可见的包括太阳、月亮、星辰、大地和水；不可见的包括神祇、魔鬼和精灵？因为他们无法否认这些存在，尽管他们可能很想否认。此外，还有其他长寿的对象，包括动物和植物。动物方面，例如鸟类中，如他们所说，有鹰、乌鸦和凤凰；陆地生物中有鹿、大象和龙；水生生物中有鲸鱼之类的海怪。植物方面，有棕榈、橡树和波斯树；树木中也有一些常青的，有人列举了十四种；另有一些只在特定季节开花，然后落叶。还有其他对象——实际上构成了绝大多数生长或生成之物——早逝而生命短暂。其中也包括人类本身，正如某部圣经论及他时所说：\BibleQuote{凡从女人生的，时日短少。}\BibleRef{约 14:1}。然而，我想他们会回答说，原子的不同组合完全可以解释如此巨大的寿命差异。因为有人认为，有些事物被原子压缩在一起，紧密交织，形成紧密结合的物体，因此极其难以分解；而另一些事物中，原子的结合或多或少较为松散和脆弱，因此或快或慢地脱离其有序结构。此外，有些物体由特定种类和共同形状的原子组成，而另一些则由不同种类、不同排列的原子组成。那么，谁是那位睿智的鉴别者，使某些原子结合，使另一些分离；并以某种方式排列一些原子以形成太阳，又以另一种方式排列另一些原子以产生月亮，并按自然适宜性和各星体的适当结构进行调整？因为那些太阳原子，以其特殊的大小和种类以及特殊的结合方式，绝不会自我缩减以产生月亮；同样，那些月亮原子的结合也绝不会发展成太阳。同样，灿烂的大角星也不会夸耀它拥有启明星的原子，昴星也不会因由猎户星的原子构成而自豪。因为保禄恰当地表达了这种区别，他说：\BibleQuote{太阳有太阳的光辉，月亮有月亮的光辉，星辰有星辰的光辉，而且星辰之间的光辉也各不相同。}\BibleRef{格前 15:41}。如果原子之间的结合是无智慧的，如同无灵魂之物的情况，那么它们必然需要一位睿智的工匠；如果它们的结合没有意志的决定而只是必然的，如同无理性之物的情况，那么必须有一位巧妙的领袖将它们聚集起来并置于其管理之下。如果它们自发地联结起来以完成自发的工作，那么必须有一位卓越的建筑师为它们分配工作并承担监督；或者必须有一位像爱秩序和纪律的将军那样对待他们，不使军队处于混乱状态，也不让一切陷入混乱，而是按顺序排列骑兵，按阵列布置重装步兵，标枪手自成一体，弓箭手分开，投石手同样如此，并将每支队伍安置在适当的位置，以便所有装备相同的人协同作战。但如果这些教师认为这个比喻只是个玩笑，因为我将非常大的物体与非常小的物体进行比较，那么我们可以转向最小的事物。\par

\Emph{然后下文接着说}：——但若没有管治者的言语、抉择和秩序加在他们身上，而凭自己行为在那汹涌激流的巨大动荡中保持正直，并在碰撞的巨响中安全渡过自己，又若像原子般相遇并同类相聚，并非如诗人所想由天主聚集，而是自然认出彼此间的亲和关系，那么，我们这里真是有一个最奇妙的原子民主：朋友欢迎并拥抱朋友，都渴望同住在一个居所；有些原子凭自己的决断把自己滚成那个巨大的发光体太阳，以形成白昼；另一些则形成许多金字塔状的星火，或许还加冕整个天空；还有一些把自己缩减成圆形，以便给以太某种坚固性，并为之拱顶，构成一个巨大的逐渐上升的光体阶梯，其目的也是让较普通原子的各种组合为自己选择定居点，并在它们之间划分天空作为住所和驻地。\par

\Emph{然后，在说了某些其他事之后，论述继续如下}：——但轻率之人甚至看不见明显的事，当然也远未察觉不明显的事。因为他们似乎甚至对天体那些有规律的升起和落下——太阳的升起和落下，连同其一切奇妙荣耀，与其他天体无异——毫无概念；他们似乎也没有恰当利用通过这些给人类提供的帮助，如白天明朗地升起供人劳作，黑夜遮蔽大地供人安息。\Quote{因为人}，经上说，\Quote{出去工作，劳作直到晚上。}他们也不考虑另一种运行，即太阳凭借其组成的那些原子，为我们确定分明的时刻、适宜的季节和规则的交替。然而，即使这样的人——他们是可怜的人，无论他们如何自以为义——可能不愿承认，却有一位伟大的主创造了太阳，并凭祂的话语赋予它运行的动力。你们这些盲人啊，你们的这些原子为你们带来冬季和雨水，好使大地为你们和地上所有生物产出食物吗？它们也带来夏季，好使你们从树上收取果实享用吗？那么，你们为什么不崇拜这些原子，向它们献祭，视它们为地上果实的守护者呢？你们实在忘恩负义，竟不把从它们那里领受的丰盛恩惠中哪怕最微薄的初果庄严地分给它们。\par

\Emph{稍作停顿后，他继续如下}：——再者，那些形成如此众多而多样的团体的星辰——这些游荡且不断自行扩散的原子所构成的星辰——借一种盟约划分了各自产业的区域，像殖民地和政府一样分配，却没有奠基人或家主的主持；它们以忠诚与和平为誓，尊重与邻星的边界法则，并避免超出最初接受的界限，仿佛在原子中享有真正君王的立法管理。然而，这些原子并不施行统治。因为什么都不是的东西怎能做这事？但请听神圣的训言：\Quote{上主的化工，都是按判断而行；从起初，在创造之时，祂就安排了它们的部分。祂永远装饰了祂的化工，它们的本原也归于它们的世代。} \BibleRef{德 16:26}\par

\Emph{然后，又过了一会儿，他继续如下}：——什么方阵曾如此完美有序地穿过平原，没有骑兵超前或落后，或阻塞道路，或让自己被阵列中的同伴抛下，就像那稳步前进、似乎持盾密集、紧凑而不间断、不受干扰、毫无阻碍的星辰队伍的行进那样？尽管由于某些偏斜和侧翼运动，它们的一些运行显得不太清晰。但无论如何，它们始终遵守指定的时期，又坚决地朝着各自升起的方位前进，仿佛那是它们精心研究的对象。因此，让这些著名的原子解剖家、这些不可分割之物的分割者、这些不可混合之物的混合者、这些领会无限的行家告诉我们，天体的这种圆周行进和路线从何而来？在这种行进中，并非任何单个原子的组合偶然意外地以这种方式旋转；而是一个巨大的圆形合唱团如此永恒地、和谐地运动，并在这些轨道上旋转。而且，这群庞大的同行者，没有队长指挥，没有意志决定，彼此互不相识，何以能如此完美和谐地保持它们的路线？先知确实把这桩事列为不可能且无法证明的事之一——即两个陌生人应一起行走。因为经上说：\Quote{二人若不相识，岂能同宿？}\par

% structure: p0022 source=h3 target=subsection reason=relative-heading-rank; base=h2

\subsection{四、从人体构造出发驳斥同一异端。}

此外，那些人不了解自己，也不了解什么是属于他们自己的。因为如果这邪恶教义的任何领袖只思考他自己是怎样的人、他来自何处，他肯定会像获得自我理解的人一样，被引导至明智的决定，并会对他的父及造物主，而不是对这些原子说：\BibleQuote{祢的手塑造了我，又再造了我。}他也会采纳这个奇妙的关于他形成的描述，正如古时一人所给予的：\BibleQuote{祢岂不是将我倾出如奶，使我凝结如乳饼？祢以皮肤和肉覆护了我，用骨与筋编织了我。祢赐予我生命和恩宠，祢的眷顾保存了我的灵魂。}因为\Person{伊壁鸠鲁}的父亲在生伊壁鸠鲁时，从他自身所发出的原子，其数量为何，其来源为何？当它们在母亲子宫内被接受时，它们如何凝聚、成形并具有形象？它们如何被推动并增长？那个微小的生殖种子如何将许多将要构成伊壁鸠鲁的原子聚集起来，将其中一些变为皮肤和肉作为覆盖，将另一些变为骨骼以保持直立和力量，又将其他变为筋腱以形成紧凑的结构？它如何构造和适配许多其他成员和部分——心、肠、感官器官，有些在内有些在外——使身体成为有生命之物？因为所有这些事物中没有一样是闲散或无用的：即使其中最卑贱的——头发、指甲或类似物——也不是如此；反而一切都有其服务要做，一切都有其贡献要做出，有些为了身体构造的健康，有些为了外表的美观。因为天主眷顾不仅关心有用的事物，也关心适时和美丽的事物。于是，头发是整个头部的保护和覆盖物，胡须是哲学家得体的装饰。因此，是天主眷顾形成了整个人体构造的一切必要部分，赋予所有成员彼此适当的连接，并从宇宙的资源中为他们量出丰富的供应。其中最主要的部分，通过个人经验的证明，甚至对未受教育者也清楚地显示出它们所附着的价值和功能：例如，头处于至高无上的位置，感官像卫兵一样围绕大脑，如同君主在城堡中；眼睛向前注视，耳朵报告信息；味觉如同贡赋的征收者；嗅觉追踪和搜索其对象；触觉操控一切置于其下的事物。\par

因此，目前我们仅以概括的方式略述全智的天主眷顾的若干作为；稍后，若天主恩准，当我们对一位被认为更有学问的人讲话时，我们将更详细地予以探讨。那么，我们有手的服务，藉此完成各类工作，实践各种技艺，这些技艺各自拥有不同的功能，却共同致力于一项普遍职能；我们有肩膀，能负重；有手指，能抓握；有肘部，能弯曲，既可向内转向身体，也可向外倾斜，从而能将物体拉向身体或推离身体。我们还有脚的服务，藉此整个受造的大地都归于我们权下，大地本身被穿越，海洋变得可航行，河流被跨越，万物之间建立了交往。腹部是食物的仓库，其各部分各得其所，合理安排，以便为自己分配适量的营养，并排出多余之物。其他一切使人身体构造得以明智地妥善管理的部分也是如此。聪明人和不聪明的人都同样享有这些用途，但对它们的理解却不相同。因为有些人将整个管理归之于他们认为真神的一位权能，他们视之为万物中至高的理智和自身最佳的恩主，相信这管理是超乎一切、本身真正神圣的智慧和能力的作为。另一些人则漫无目标地将这整个极其美妙的构造归之于偶然和巧合。此外，还有一些医生，他们对这一切进行了更有效的检查，极其精确地研究了内部器官的排列，对其探究的结果感到惊异，从而将自然本身神化。这些人的观念我们日后将尽可能加以审视，尽管可能只触及表面。与此同时，为了一般而概括地处理这个问题，请问是谁建造了我们这整个居所，使其如此崇高、挺直、优雅、敏锐、可动、活跃，且适于一切？是如他们所言的那无理性的原子群吗？不，这些原子即使结合，也不能捏塑一个泥像，不能雕琢和打磨一尊石像，也不能铸造和完成一尊银像或金像；相反，制造这些物品的人们已发现了适于这些操作的艺术和技艺。如果连这些仿制品和模型都不能在没有智慧的情况下做出，那么这些仿制品所依据的真实原型又如何能自发地产生？灵魂、理智和理性与哲学家一同出生，它们从何而来？难道他从那些没有灵魂、理智和理性的原子中收集了这些？这些原子中的每一个是否都向他灌输了某种适当的观念和概念？我们是否要认为，人的智慧是由这些原子组合而成，正如赫西俄德的神话所说潘多拉是由众神塑造的？那么，希腊人就不得不放弃谈论各种诗歌、音乐、天文、几何以及其他一切艺术和科学是神的发明和教导，而必须承认这些原子是唯一精通一切题材的、拥有技艺和智慧的缪斯。因为伊壁鸠鲁用原子构造的\WorkTitle{神谱}，确实是某种外在于无限有序世界的东西，并在无限的无序中找到避难所。\par

% structure: p0025 source=h3 target=subsection reason=relative-heading-rank; base=h2

\subsection{五、为天主工作不是一件辛苦和疲劳的事。}

如今，工作、管理、行善、照顾及类似的行为，对于懒惰、愚蠢、软弱和邪恶的人来说，也许是艰巨的任务——而\Person{伊壁鸠鲁}确实把自己算在其中，当他提出关于诸神的这样一些观念时。但对于认真、有力、明智和谨慎的人——而哲学家应当如此，何况诸神呢！——这些东西不仅不是令人不愉快和厌烦的，反而总是最令人愉快、最受欢迎的。对于这种性格的人，在行善上的疏忽和拖延是一种耻辱，正如诗人用以下劝诫的话所告诫他们的：——\par

\Quote{不要将任何事情推迟到明日}\par

\Quote{然后他又补充了以下威胁的话：——}\par

\Quote{懒惰的拖延者永远与痛苦搏斗。}先知也以更庄重的方式教训我们，并宣称按德行标准行事确实配得上天主，而不留意这些的人是有祸的：\BibleQuote{为主工作疏忽的人，是可咒骂的。}此外，那些不精通任何技艺、不能完美从事的人，在初次尝试时会感到厌倦，因为经验新奇且缺乏实践。但那些已经进步的人，更不用说完全训练有素的人，会轻松成功地实现他们的劳动目标，并在工作中获得极大乐趣，宁愿这样完成他们习惯的追求并完美实现努力目标，也不愿获得所有被人们认为有利的东西。是的，据说\Person{德谟克里特}自己断言，他宁愿发现一个真正的原因，也不愿取得\Place{波斯}王国。这是一个仅对事物原因有虚妄无根据概念的人的声明，因为他从一个没有根据的原则和错误假设出发，没有辨别自然事物构成中的真正根源和共同的\Emph{必然律}，并把对简单、无理智和随机发生之事的把握视为最大智慧，并设立机运为一切事物甚至神性事物的女皇，努力证明一切事物都因机运的决定而发生，但同时将机运排除在人类生活之外，并指责那些崇拜它的人愚昧。无论如何，在他的\WorkTitle{格言集}开头，他这样说：\Quote{人们制造了一个机运的形象，作为他们自己无知的遮盖。因为理智和机运本质上彼此对立。人们坚持认为这个对理智的最大敌手是理智的主宰。是的，他们完全颠覆并废除了理智，而将机运置于其位。因为他们不赞美理智是幸运的，反而颂扬机运是最聪明的。}此外，那些关注有益于生活之事的人，特别乐于从事有利于同族利益的事，并寻求因承担公共利益的劳动而获得的赞美和荣耀的回报；有些人作为方式方法的提供者，有些人作为官员，有些人作为医生，有些人作为政治家；甚至哲学家也为教育人类的努力而自豪。那么，\Person{伊壁鸠鲁}或\Person{德谟克里特}敢断言，在哲学思考的努力中，他们只给自己带来痛苦吗？不，他们会认为这是无与伦比的头脑愉悦。因为即使他们坚持善即快乐的观点，他们也会羞于否认哲学思考对他们来说是更大的快乐。至于诸神，他们的诗人歌颂为\Quote{恩赐的赐予者}，这些哲学家却嘲弄地这样歌颂他们：\Quote{诸神既不是任何好事的赐予者，也不是分享者。}确实，他们能如何证明诸神存在呢？当他们既感知不到他们的临在，也看不出他们做任何事时，这方面他们就像那些崇拜太阳、月亮和星辰并称它们为\Emph{诸神}（从其\Emph{运行}得名）的人一样；此外，他们不赋予诸神任何功能或运作能力，从而不因\Emph{构成}（即\Emph{创造}对象）而称它们为神（因为据此，一切事物的唯一创造者和运作者必定是天主）；最后，他们没有提出诸神对人类有任何管理、审判或恩惠，使我们因恐惧或崇敬而崇拜他们？那么，\Person{伊壁鸠鲁}真的能看透这个世界并跨越天界吗？或者他通过只有他自己知道的秘密大门出去，从而在虚空中看到了诸神？认为他们完全幸福，然后自己成为这种快乐的热切追求者，以及他们在虚空中所追求的生活的热心学生，现在他呼吁所有人参与这种幸福，并敦促他们这样使自己像诸神一样，为他们准备真正的\Emph{会饮}——不是天堂或\Place{奥林帕斯}，如诗人所虚构的，而是纯粹的虚空，并摆上原子的神馔，用同样的东西调制的仙酒祝酒？然而，在与我们无关的事物上，他在书中引入了无数誓言和庄严的断言，不断地向朱庇特发誓肯定和否定，并让他遇到的人和他讨论他的学说的人也向诸神发誓，当然不是他自己害怕他们，或害怕伪善，而是他认为这一切都是虚妄、虚假、空洞和不可理解的，只是作为他话语的一种伴奏，就像他可能清喉咙、吐痰、扭曲脸或移动手一样。在他看来，整个诸神命名之事是完全荒唐空洞的假装。但这也是一个非常明显的事实：他害怕\Person{苏格拉底}之死后的雅典人，并希望防止被人当作他真正的样子——无神论者——这个狡猾的骗子为他们发明了一些空洞的虚幻诸神的影子。但他肯定从未用真正的理智之眼仰望天空，以便听到来自高处的清晰声音——另一个专注的观察者确实听到了，并作证说：\BibleQuote{诸天宣扬天主的荣耀，穹苍显示他手的化工。}他也肯定没有用应有的反思俯视世界表面！因为那样他会学到\BibleQuote{大地充满主的良善}和\BibleQuote{大地和其中的万物属于主}；并且，正如我们读到：\BibleQuote{此后，主垂顾大地，以他的祝福充满它。他用各种生物覆盖了它的表面。}如果这些人尚未完全盲目，让他们看看生物的丰富多样性：陆地动物、飞鸟和水生动物；然后他们就会明白主在审判万物时所作的宣告是真确的：\BibleQuote{一切事物，按他的命令，都显得美好。}\par

% structure: p0030 source=h2 target=section reason=relative-heading-rank; base=h2

\section{三、摘自\WorkTitle{驳撒贝里乌斯}。论及物质非受生的观念。}

那些认为物质非受生的人，当然不能被视为虔诚。他们固然承认物质在天主的手下，就其安排与规整而言，是被支配的；因为他们承认物质本质上是被动与可塑的，容易随从天主所施加的改变。然而，他们必须向我们清楚说明，相似与不相似怎能同时共存于天主与物质中。因为如此一来，就必须认为有一个高于两者的存在，这种思想不能合理地用于天主。因为如果两者都有这种非受生性的欠缺（这被视为相似之处），而两者又有此差异点，那么这种差异从何而来？如果天主是非受生的，并且这种非受生的欠缺可以说是祂的本质，那么物质就不能是非受生的；因为天主与物质并非同一。但如果两者各自独立存在——即天主与物质——并且非受生的欠缺也属于两者，那么显然有一个与两者不同、比两者更古老更高的存在。但两者相反构造的差异完全推翻了它们能够平等共存的想法，更推翻了两者中的一方——即物质——能够独立存在的想法。因为那样他们将不得不解释：虽然两者都被假定为非受生，但天主却是不可受难、不变、不动、充满动力的，而物质则相反：可受影响、可变、易变、可改。那么，这些属性怎能和谐共存与结合呢？难道是天主使自己适应物质的本性，从而巧妙地加工它？但假设天主像人一样加工金子，或砍凿石头，或从事其他使不同物质能接受形状样式的技艺，那是荒谬的。然而，另一方面，如果祂按自己的旨意，根据自己智慧的指示，将祂自己运作所产生的丰富多样的形式印在物质上，那么我们的这个论述就是美好而正确的，并且进一步确立了非受生的天主是宇宙万物的\OriginalTerm{本体}{hypostasis}（生命和基础）这一立场。因为这一非受生的事实与祂存在的得当方式相结合。要驳斥这些教师，还有很多可说的，但这并非我们当前的任务。如果把他们与最不虔敬的多神论者并列，后者在言辞上似乎更为虔敬。\par

% structure: p0032 source=h2 target=section reason=relative-heading-rank; base=h2

\section{四、致罗马主教狄约尼削书}

% structure: p0033 source=h3 target=subsection reason=relative-heading-rank; base=h2

\subsection{摘自第一卷。}

1. 天主为父，决无一时不是父。\par

2. 确实，并非好像天主没有生出这些，后来才生圣子，而是因为圣子的存在并非来自自己，而是来自圣父。\par

此后数语，他论及圣子说：——\par

3. 既然祂是永恒之光的辉耀，祂本身也绝对是永恒的。因为光总是存在的，显然它的辉耀也存在，因为光的存在是通过它发光而感知的，而光不可能不发光。让我们再次举例说明。如果太阳存在，也就有白昼；如果这些都不显现，太阳就不可能存在。那么，如果太阳是永恒的，白昼就永无终结；但如今，情况并非如此，白昼与太阳一同开始，也一同结束。但天主是永恒之光，既无开始，也永不消失。因此，永恒的辉耀在祂面前照耀，并与祂同存，因为它无始而生，永远在祂面前照耀；而祂就是那智慧，它说：\BibleQuote{我是祂所喜悦的，我日日作祂的喜乐，常在祂面前。}\par

稍后，他又从同一点继续论述说：——\par

4. 因此，既然圣父是永恒的，圣子也是永恒的，光之光。因为哪里有生者，哪里就有所生者。如果没有所生者，祂怎能、又如何是生者？但两者都存在，且永远存在。既然天主是光，基督就是辉耀。既然祂是神——因为祂说：\BibleQuote{天主是神}——同样，基督也被称为气息，因为祂说：\BibleQuote{祂是天主威能的气息。}\par

他又说：——\par

5. 再者，圣子独自常与圣父同存，并为那自存者所充满，祂本身也存在，因为祂出自圣父。\par

% structure: p0042 source=h3 target=subsection reason=relative-heading-rank; base=h2

\subsection{摘自同一卷书。}

6. 但当我说及受造之物与某些需要思考的作品时，我匆忙提出了一些或许不太恰当的比喻，例如我说，植物与农夫并不相同，船与造船者也不相同。但随后我更多地停留在更合宜且更符合事物本性的事物上，我以许多言语、通过各种周密考虑的论证，阐述了哪些事物更为真实。此外，我在另一封信中已将这些事向您阐明。在这些事情上，我也证明了他们对我提出的指控的虚假性，即他们声称我不坚持基督与天主\OriginalTerm{同质}{consubstantial}。因为尽管我说过，我从未在\WorkTitle{圣经}中发现或读到这个词，但我随即添加的其他推理与这一观点毫不相悖，因为我引用了人类的子孙作为比喻，这无疑与生育者属于同一种类。并且我说过，父母与子女绝对区别的唯一事实就是他们自己不是其子女，否则便必然既没有父母也没有子女。但是，正如我之前所说，由于目前的情况，我手头没有那封信；否则我会将我当时写下的原话寄给您，是的，还有整封信的副本。如果将来有机会，我会寄给您。我还记得，我添加了许多来自彼此相似之物的比喻。因为我说过，植物，无论是从种子还是从根生长出来，都与它发芽的来源不同，尽管它绝对具有相同的本性；同样，从泉源流出的河流具有另一种形式和名称：因为泉源不被称为河流，河流也不被称为泉源，但它们是两样事物，而且泉源仿佛是父亲，而河流是来自泉源的水。但他们假装看不见这些以及相似于这些的文字，仿佛他们是盲人；但他们试图用过于轻率的措辞从远处攻击我，就像扔石头一样，没有注意到，对于他们所不了解且需要解释才能理解的事物，不仅是不相似的比喻，甚至是相反的比喻，也常常会阐明它们。\par

% structure: p0044 source=h3 target=subsection reason=relative-heading-rank; base=h2

\subsection{出自同一第一卷。}

7. 上面说过，天主是一切善事的泉源，而圣子被称为从祂流出的河流；因为圣言是心灵的一种流出，并且——按人的说法——由口从心中发出。但由舌头发出的心灵与存在于心中的圣言不同。因为后者，在发出前者之后，仍然保持原来的样子；而被发出的心灵则飞离，到处流转，因此两者彼此内在于对方，尽管一者出于另一者，它们是一体，尽管是两者。同样，圣父与圣子被称为一体，并且彼此内在于对方。\par

% structure: p0046 source=h3 target=subsection reason=relative-heading-rank; base=h2

\subsection{出自第二卷。}

8. 我所说出的各个名字既不能彼此分离，也不能分开。我提到了圣父，在我提及圣子之前，我已经在圣父中暗示了祂。我加上了圣子，而圣父，即使我之前未予指名，也已完全包含在圣子内。我加上了圣神，但同时，我以那名字传达了祂由谁和由谁发出。但他们不知道，圣父，作为\Emph{圣父}，不能与圣子分离，因为那个名字是明显的合一与联系的依据；圣子也不能与圣父分离，因为\Emph{圣父}这个词指示了他们之间的关联。此外，还有明显可见的圣神，祂既不能与派遣祂的那一位分离，也不能与带来祂的那一位分离。那么，我使用这些名字，又怎么会认为它们是完全分开、彼此分离的呢？\par

数语之后，他补充道：——\par

9. 确实，我们将不可分割的统一体扩展为天主圣三，又将不可减少的天主圣三收缩为统一体。\par

% structure: p0050 source=h3 target=subsection reason=relative-heading-rank; base=h2

\subsection{出自同一第二卷。}

10. 但如果任何吹毛求疵者，因我曾说过天主是万物的制作者和创造者，就认为我也说过祂是基督的创造者，那么请他注意，我首先称祂为圣父，这个词同时表达了圣子。因为在我称圣父为创造者之后，我补充说：如果生育者被正确理解为父亲（因为我们将在接下来的内容中考虑\Emph{圣父}一词的广度），祂也不是祂所创造之物的父亲。制作者也不是父亲，如果只有塑造者才被称为制作者。因为在希腊人中，智者被称为其著作的制作者。宗徒也说：\Quote{行律法的人（即制造者）。}此外，对于内心的事，如美德与恶习，人们被称为行事者（即，制造者）；按此方式，天主说：\Quote{我期望它施行判断，它却行了不义。}\par

11. 这话也不该如此被责备；因为他说，他使用\Quote{制造者}之名是由于圣言所取的人性，这人性确实是受造的。但如果有人怀疑那是指圣言而说的，那么这也应当不争辩地听。因为正如我不认为圣言是受造物，我也不说天主是它的制造者，而是它的圣父。然而，如果有些时候，在谈论圣子时，我可能偶然说过天主是祂的制造者，即使这种说法也不是没有辩护的。因为希腊人中的智者称自己为著作的制造者，虽然他们也是其著作的父亲。此外，神圣的\WorkTitle{圣经}称我们为那些发自内心的行动的制造者，当它称我们为判断和正义之律法的实行者时。\par

% structure: p0053 source=h3 target=subsection reason=relative-heading-rank; base=h2

\subsection{出自同一第二卷。}

12. \Emph{在起初已有圣言}。但那并不是产生圣言的圣言。因为圣言与天主同在。主是智慧；因此并非智慧产生智慧；因为祂说：\Quote{我是那}，祂在其中喜悦，基督是真理；但祂说：\Quote{有福的}，\Quote{是真理的天主。}\par

% structure: p0055 source=h3 target=subsection reason=relative-heading-rank; base=h2

\subsection{选自第三卷。}

13. 生命由生命所生，如同河流从泉源涌出，灿烂的光从不灭之光点燃。\par

% structure: p0057 source=h3 target=subsection reason=relative-heading-rank; base=h2

\subsection{选自第四卷。}

14. 正如我们的心智从自身发出言语——如先知所说：\Quote{我的心里涌出了优美的言语}——二者彼此有别，各自保持独特的位置和区分；前者存留并运行在心中，后者位于舌头和口唇。然而它们并不相离，也不相缺；心智没有言语，言语没有心智；而是心智产生言语，并显现在言语中，言语展示那产生它的心智。心智仿佛是内在的言语，言语是流露的心智。心智进入言语，言语将心智传递给周围的听者；如此，心智借着言语进入听者的灵魂，与言语同时进入。的确，心智仿佛是言语的父亲，存在于自身；言语仿佛是心智的儿子，不能先于心智或没有心智而产生，而是与心智共存，从中获得种子和起源。同样，全能圣父和宇宙心智在万有之先就有了圣子、圣言和话语，作为祂的解释者和使者。\par

% structure: p0059 source=h3 target=subsection reason=relative-heading-rank; base=h2

\subsection{约在论文中部。}

15. 如果因有三个\OriginalTerm{位格}{hypostases}，他们就说它们是被分割的，那么无论他们愿意与否，都有三个，否则就让他们完全取消神圣的天主圣三。\par

又说：\par

因为为此缘故，在统一之后，还有最神圣的天主圣三。\par

% structure: p0063 source=h3 target=subsection reason=relative-heading-rank; base=h2

\subsection{整篇论文的结论。}

16. 依照这一切，我们从先于我们生活的长老领受了形式和规则，我们也以与他们一致的声音，向你们表达感谢，并为现在所写的信函。愿光荣与权能归于天主圣父、及祂的圣子我们的主耶稣基督，同着圣神，直到永远。阿们。\par

% structure: p0065 source=h2 target=section reason=relative-heading-rank; base=h2

\section{五、致巴西里德主教的信。}

% structure: p0066 source=h3 target=subsection reason=relative-heading-rank; base=h2

\subsection{第一规条。}

狄约尼削致我亲爱的儿子和兄弟、在圣事中与我同工的、天主的顺服仆人巴西里德，在主内问安。\par

最忠信且成全的吾子，你曾致书于我，询问在五旬节日，应当于何时结束守斋。因为你说，有些弟兄认为当在鸡鸣时，另有些则认为当在黄昏时。在罗马的弟兄们，如他们所说，等候鸡鸣；而关于此地的弟兄，你告知我们说，他们愿意更早结束。因此你急切地盼望能够准确无误地确定那个时辰，这确实是一件困难且不确定的事。然而，众人必会欣然承认：从吾主复活之日起，那些直到那时仍以斋戒谦卑自己灵魂的人，应当即刻开始他们的节期喜乐与欢欣。但你在致我的信函中，以对圣经明智的理解清楚地表明，圣经中似乎并未就祂复活的时辰提供十分精确的记载。因为福音作者们对不同时间来到坟墓的人有不同描述，但都宣布他们发现主已经复活了。正如玛窦所说：\\BibleQuote{安息日将尽}；若望记载：\\BibleQuote{清晨，天还黑}；路加记述：\\BibleQuote{一周的第一天，大清早}；马尔谷告诉我们：\\BibleQuote{清早，太阳升起时}。如此，没有一人明确地告诉我们祂复活的确切时辰。然而，众人承认，那些在安息日将尽、一周的第一天破晓时来到坟墓的人，发现祂已不再躺在那里。我们不要以为福音作者们彼此不合或矛盾。即使关于我们所探究的问题似乎有些微困难，倘若他们一致同意世界之光、我们的主是在那一夜复活，只是对时辰有不同看法，我们仍宜以智慧而忠实的心智调和他们的陈述。玛窦的叙述如下：\\BibleQuote{安息日将尽，一周的第一天天亮时，玛利亚玛达肋纳和另一位玛利亚来看坟墓。看，发生了大地震，因为主的天使从天降下，前来把石头滚开，坐在上面。他的容貌好像闪电，他的衣服洁白如雪；看守的人由于怕他，战抖起来，好像死人。天使对妇女们说：\\InnerQuote{不要害怕！我知道你们寻找被钉死的耶稣。祂不在这里，因为祂已经照祂所说的话复活了。}}如今，有些人会按照惯常用法认为\\Quote{将尽}一词表示安息日的\\Emph{晚上}；另有些人，对事实有更深的认知，会说它不指示\\Emph{晚上}，而是\\Emph{深夜的末时}，因为\\Quote{将尽}一词表示缓慢而漫长的时间。而且因为他提到\\Quote{夜}，而非\\Quote{晚上}，他加上了\\Quote{一周的第一天天亮时}这句话。这里的人尚未如他人所说\\Quote{带着香料}前来，而是\\Quote{来看坟墓}；他们发现了地震，天使坐在石头上，并听到天使宣布：\\Quote{祂不在这里，祂已经复活了。}若望的见证也与此一致。他说：\\BibleQuote{一周的第一天，清晨，天还黑，玛利亚玛达肋纳来到坟墓，看见石头已从坟墓挪开了。}只是按照这\\Quote{天还黑}，她来得较早。路加说：\\BibleQuote{她们在安息日，按照诫命安息了。在一周的第一天，大清早，她们来到坟墓，带着预备好的香料，发现石头已从坟墓滚开了。}这\\Quote{大清早}很可能指示一周的第一天破晓时分；如此，安息日本身完全过去，随后的整夜也过去，另一天开始，她们带着香料和没药前来，那时才显明祂早已复活。马尔谷跟随此说：\\BibleQuote{她们买了香料，好来傅抹祂。一周的第一天，大清早，她们来到坟墓，太阳升起时。}因为这位福音作者也用\\Quote{大清早}，与前者所用\\Quote{大清早}相同；他又加了\\Quote{太阳升起时}。如此，她们出发，起初在\\Quote{大清早}或（如马尔谷所说）\\Quote{大清早}上路；但在路上，以及在坟墓逗留期间，她们一直待到太阳升起。然后那个穿白衣的少年对他们说：\\BibleQuote{祂已经复活了，不在这里了。}既然情形如此，我们向那些寻求确切时辰（甚至半时、一刻）以开始庆祝主从死者中复活的人，作如下陈述与解释。那些过于急躁，甚至在午夜前就结束的人，我们斥之为懈怠放纵，仿佛因仓促而几乎中断了行程，因为智者的言语说：\\Quote{生命中小事亦非小事。}那些坚持很长时间，甚至持续到第四更（也是我们的救主显现自己在海上行走于那时在深海之人面前的时候）的人，我们视之为高贵而勤苦的门徒。至于那些在行程中因感动或力所能及而停下来歇息的人，我们不必过于严厉：因为并非所有人都同样或一致地度过六天斋戒；有些人整日守斋，完全不吃东西；另有些人只守两天，或三天，或四天，甚至有人未守一天。对于那些辛劳地持守长期斋戒，而后精疲力竭、几乎无法支撑的人，若他们进食稍早，宜予宽恕。但若有任何人不仅拒绝如此长期的斋戒，而且起初完全不愿守斋，反而在前四日尽情享乐，及至最后两日——即预备日和安息日——才严格守斋，且只在这两日守，并认为若能坚持到早晨便行了伟大光辉之事，我无法认为他们与那些提前几日就操练的人平等地度过了时间。这是我根据对问题的理解，以书面形式就这些事向你提供的劝告。\par

% structure: p0069 source=h3 target=subsection reason=relative-heading-rank; base=h2

\subsection{第二会规。}

关于妇女在经期是否适宜进入天主的圣殿这一问题，我认为这是一个多余的询问。因为我相信，如果她们是虔诚的信女，她们不会鲁莽到在这种情形下接近祭台或触摸主的身体和血。诚然，那患血漏十二年的妇人并没有触摸\Emph{主自己}，而只是摸了祂的衣边以求痊愈。因为无论一个人处于何种境况，祈祷、记念主、并祈求帮助，都是无可指摘的操练。但那在灵魂和身体上都不完全纯洁的人，应被阻止接近至圣所。\par

% structure: p0071 source=h3 target=subsection reason=relative-heading-rank; base=h2

\subsection{第三会规。}

此外，凡有能力的且年长者，应在这些事上做自己的判断。因为他们已从保禄的书信听到，应当彼此同意分房，以便专心祈祷，然后再同房。\par

% structure: p0073 source=h3 target=subsection reason=relative-heading-rank; base=h2

\subsection{第四会规。}

至于那些在夜间被不自主的遗精所困扰的人，让他们随从自己良心的见证，并判断自己在这事上是否有疑虑。关于食物，宗徒告诉我们：\Quote{若疑惑而吃，就被定罪。}因此，在这些事上，凡接近天主的人都应具有良善的良心和适当的自信，凭着自己的判断。确实，为表示你对我们的尊重（因为你并非无知，亲爱的），你向我们提出了这些问题，使我们同心合意，正如我们与你一样同心同神。而我，则公开陈述了我的意见，并非作为导师，而只是如我们彼此间以坦诚相商。我最有智慧的儿子，当你审查我这意见后，请写信告知我你对这些事的看法，并让我知道任何你认为公正和可取的观点，以及你是否赞同我的判断。愿我亲爱的儿子平安地事奉主，这是我的祈祷。\par

\SourceRecord{Translated by S.D.F. Salmond.；Ante-Nicene Fathers；Vol. 6.；Edited by Alexander Roberts, James Donaldson, and A. Cleveland Coxe.；Buffalo, NY: Christian Literature Publishing Co.,；1886.；<http://www.newadvent.org/fathers/0612.htm>.}
